% ---------------------------------------------------------------------------
% Appendix: API reference
% ---------------------------------------------------------------------------
\chapter{API Reference}\label{app:api}
% Long monospaced identifiers leave little room for line breaking; allow some
% extra inter-word stretch in this chapter (the group keeps it local).
\begingroup\setlength{\emergencystretch}{3em}% chapter-local

This appendix lists the public classes and functions of version 1.0 with their
signatures, grouped by module. Signatures omit type annotations; defaults are
shown as in the source code. The chapters give the semantics, and the
docstrings (\code{help(obj)} in Python) give full details. Names that start
with an underscore and modules not listed here are internal and may change
without notice.

\section{\texorpdfstring{\pkg{env\_lib}}{env\_lib}}\label{sec:api-envlib}

\subsection{Package functions}

\begin{pycode}
env_lib.make(env_id, **kwargs) -> gym.Env
env_lib.list_envs() -> list[str]
env_lib.register_envs()
env_lib.__version__                       # "1.0.0"
env_lib.ResetNeededError                  # exception class
\end{pycode}

\code{make} creates a registered environment (a thin wrapper around
\code{gymnasium.make} that suppresses the ``out of date'' warning of
\code{-v0} ids); keyword arguments go to the environment constructor.
\code{list\_envs} returns the registered ids in the order of
Table~\ref{tab:api-ids}. \code{register\_envs} registers all ids with
Gymnasium; it runs automatically on \code{import env\_lib} and is idempotent.
\code{ResetNeededError} is raised by every environment when \code{step()} or
\code{render()} is called before \code{reset()}; it subclasses both
\code{RuntimeError} and \code{gymnasium.error.ResetNeeded}.
The environment classes below are exported lazily by \pkg{env\_lib}, so
optional dependencies are imported only when a class is used.

\begin{table}[htbp]
  \centering
  \small
  \caption{Registered environment ids and their preset arguments.}
  \label{tab:api-ids}
  \begin{tabularx}{\textwidth}{@{}>{\raggedright\arraybackslash}p{6.1cm}>{\raggedright\arraybackslash}X@{}}
    \toprule
    Id & Class and preset arguments \\
    \midrule
    \envid{KuramotoOscillator-v0} & \code{KuramotoOscillatorEnv()} (10 oscillators) \\
    \envid{KuramotoOscillator-v1} & \code{n\_oscillators=5} \\
    \envid{KuramotoOscillator-Constant-v0} & \code{n\_oscillators=6}, \code{coupling\_mode="constant"} \\
    \envid{KuramotoOscillator-\allowbreak FreqSync-Constant-v0} & as above, \code{reward\_type="frequency\_synchronization"}, \code{target\_frequency=2.0} \\
    \envid{KuramotoOscillatorTorch-v0} & \code{KuramotoOscillatorEnvTorch()} (10 oscillators) \\
    \envid{KuramotoOscillatorTorch-v1} & \code{n\_oscillators=8}, \code{n\_agents=4}, \code{device="cpu"} \\
    \envid{KuramotoOscillatorTorch-v2} & \code{n\_oscillators=10}, \code{device="auto"}, \code{integration\_method="rk4"} \\
    \envid{KuramotoOscillatorTorch-\allowbreak Constant-v0} & \code{n\_oscillators=6}, \code{coupling\_mode="constant"} \\
    \envid{KuramotoOscillatorTorch-\allowbreak FreqSync-Constant-v0} & as above, frequency synchronisation, \code{target\_frequency=2.0} \\
    \envid{LineMsg-v0} & \code{LineMsgEnv()} \\
    \envid{WirelessComm-v0} & \code{WirelessCommEnv()} ($6 \times 6$ grid) \\
    \envid{WirelessComm-v1} & \code{grid\_x=4}, \code{grid\_y=4} \\
    \envid{Pistonball-v0} & \code{PistonballEnv()} \\
    \envid{Consensus-v0} & \code{ConsensusEnv()} \\
    \envid{Formation-v0} & \code{ConsensusEnv(task="formation")} \\
    \envid{AJLATT-v0} & \code{AJLATTEnv()} (Gymnasium's passive environment checker disabled) \\
    \bottomrule
  \end{tabularx}
\end{table}

\subsection{Common environment interface}

Every environment class implements the Gymnasium methods

\begin{pycode}
reset(*, seed=None, options=None) -> tuple[observation, dict]
step(action) -> tuple[observation, reward, terminated, truncated, dict]
render() -> np.ndarray | None        # (H, W, 3) uint8 in "rgb_array" mode
close()                              # idempotent
\end{pycode}

and has the attributes \code{action\_space}, \code{observation\_space},
\code{render\_mode}, \code{metadata} and \code{np\_random}. The \code{info}
dictionary always contains \code{"agent\_rewards"}, a \code{float64} array of
shape \code{(n\_agents,)}. Chapter~\ref{chap:envlib} describes the conventions.

\subsection{\texorpdfstring{\code{KuramotoOscillatorEnv} and \code{KuramotoOscillatorEnvTorch}}{KuramotoOscillatorEnv and KuramotoOscillatorEnvTorch}}

Module \pkg{env\_lib.kos\_env}; Chapter~\ref{chap:kuramoto}.

\begin{pycode}
KuramotoOscillatorEnv(n_oscillators=10, n_agents=1, dt=0.01, max_steps=1000,
                      coupling_range=(0.0, 5.0), control_input_range=(-1.0, 1.0),
                      natural_freq_range=(0.5, 2.0), device='cpu',
                      render_mode=None, integration_method='euler',
                      reward_type='order_parameter', noise_std=0.0,
                      topology='fully_connected', adj_matrix=None,
                      target_frequency=1.0, coupling_mode='dynamic',
                      constant_coupling_matrix=None, coupling_strength=1.0,
                      normalize_coupling=False, topology_seed=42,
                      sync_threshold=0.99, sync_bonus=10.0)

KuramotoOscillatorEnvTorch(...same arguments..., device='cpu', render_agent=0)
\end{pycode}

\code{KuramotoOscillatorEnv} simulates one network with NumPy.
\code{KuramotoOscillatorEnvTorch} simulates \code{n\_agents} independent
systems in a batch on \code{device} (\code{"cpu"}, \code{"cuda"},
\code{"cuda:<k>"}, a \code{torch.device} or \code{"auto"});
\code{render\_agent} selects the rendered system. Additional methods of the
PyTorch class:

\begin{center}
  \small
  \begin{tabularx}{\textwidth}{@{}>{\raggedright\arraybackslash}p{6.9cm}>{\raggedright\arraybackslash}X@{}}
    \toprule
    Method & Description \\
    \midrule
    \code{get\_batch\_observations()} & Tensor of the observations of all systems, \code{(n\_agents, obs\_dim)}, on the device. \\
    \code{get\_batch\_rewards(dphases\_dt=None, include\_bonus=False)} & Per-system rewards of the current state, \code{(n\_agents,)}. \\
    \bottomrule
  \end{tabularx}
\end{center}

\code{env\_lib.kos\_env.register()} is deprecated; registration happens on
import.

\subsection{\texorpdfstring{\code{LineMsgEnv} and \code{WirelessCommEnv}}{LineMsgEnv and WirelessCommEnv}}

Modules \pkg{env\_lib.linemsg\_env} and \pkg{env\_lib.wireless\_comm\_env};
Chapter~\ref{chap:networked}.

\begin{pycode}
LineMsgEnv(num_agents=10, n_obs_neighbors=1, max_iter=50, render_mode=None,
           action_space_type='discrete')

WirelessCommEnv(grid_x=6, grid_y=6, ddl=2, packet_arrival_probability=0.8,
                success_transmission_probability=0.8, n_obs_neighbors=1,
                max_iter=50, render_mode=None)
\end{pycode}

Both support the render modes \code{"human"}, \code{"rgb\_array"} and
\code{"ansi"} (text). \code{seed(seed=None)} is kept for backward
compatibility and deprecated in favour of \code{reset(seed=...)}.
\code{WirelessCommEnv} also exposes \code{access\_point\_mapping(i, j,
agent\_action)}, which returns the access point \code{(row, column)} used by
agent \code{(i, j)} (or \code{(None, None)}), and
\code{check\_transmission\_fail(ap\_x, ap\_y, access\_point\_profile)}, which
is \code{True} unless exactly one agent uses the access point.
\code{env\_lib.wireless\_comm\_env.register()} is deprecated.

\subsection{\texorpdfstring{\code{PistonballEnv}}{PistonballEnv}}

Module \pkg{env\_lib.pistonball\_env}; Chapter~\ref{chap:pistonball}. Requires
the \code{pistonball} extra.

\begin{pycode}
PistonballEnv(n_pistons=20, time_penalty=-0.1, continuous=True, random_drop=True,
              random_rotate=True, ball_mass=0.75, ball_friction=0.3,
              ball_elasticity=1.5, max_cycles=125, render_mode=None,
              movement_penalty=0.0, movement_penalty_threshold=0.01, kappa=1,
              terminated_condition=True, leftmost_piston_reward=0.0,
              termination_reward=0.5)

ManualPolicy(env, agent_id=0, show_obs=False)
\end{pycode}

\begin{center}
  \small
  \begin{tabularx}{\textwidth}{@{}>{\raggedright\arraybackslash}p{6.9cm}>{\raggedright\arraybackslash}X@{}}
    \toprule
    Member & Description \\
    \midrule
    \code{observable\_mask(ball\_x=None)} & Boolean \code{(n\_pistons,)} mask of the pistons within \code{kappa} of the piston under the ball. \\
    \code{get\_local\_reward\_for\_piston(i, prev\_x, curr\_x)} & Reference implementation of one piston's local reward. \\
    \code{get\_local\_reward(prev\_position, curr\_position)} & The ball term of the local reward. \\
    \code{renderer} & The \code{PistonballRenderer}, or \code{None} before the first \code{render()}. \\
    \code{seed(seed=None)} & Deprecated; use \code{reset(seed=...)}. \\
    \code{draw()}, \code{draw\_background()}, \code{draw\_pistons()}, \code{enable\_render()} & Deprecated no-ops kept for compatibility; use \code{render()}. \\
    \bottomrule
  \end{tabularx}
\end{center}

\code{ManualPolicy} is a keyboard controller for \code{"human"} mode: calling
it returns a joint action in which the selected piston follows the keys
(W/S or Up/Down move it, A/D or Left/Right select another piston); its
\code{running} attribute becomes \code{False} when the user presses Escape
or closes the window. The module also exports \code{PistonballLayout} (screen
geometry) and \code{PistonballRenderer} (the pygame renderer).

\subsection{\texorpdfstring{\code{ConsensusEnv}}{ConsensusEnv}}

Module \pkg{env\_lib.consensus\_env}; Chapter~\ref{chap:consensus}.

\begin{pycode}
ConsensusEnv(*, n_agents=8, task='consensus', topology='ring', dynamics='single',
             dt=0.1, max_steps=200, arena_size=10.0, max_control=1.0,
             control_cost=0.01, noise_std=0.0, tolerance=0.05, success_bonus=10.0,
             formation_shape='circle', formation_radius=3.0, comm_radius=4.0,
             edge_probability=0.3, graph_seed=None, max_neighbors=None,
             damping=0.5, render_mode=None)
\end{pycode}

\begin{center}
  \small
  \begin{tabularx}{\textwidth}{@{}>{\raggedright\arraybackslash}p{6.9cm}>{\raggedright\arraybackslash}X@{}}
    \toprule
    Member & Description \\
    \midrule
    \code{laplacian\_policy(gain=1.0, velocity\_gain=None)} & Distributed feedback $u = -\text{gain}\,L(x - d)$ (plus velocity feedback for double integrators), clipped to the action bounds. \\
    \code{positions}, \code{velocities} & Copies of the current states, \code{(n\_agents, 2)}. \\
    \code{formation\_offsets} & Offsets $d_i$, \code{(n\_agents, 2)} (zeros for consensus). \\
    \code{adjacency}, \code{laplacian} & Current communication graph and its Laplacian $L = D - A$. \\
    \code{algebraic\_connectivity} & Fiedler value of the current Laplacian. \\
    \code{task\_error}, \code{step\_count} & Current task error and step within the episode. \\
    \code{observation\_layout} & Column slices of one observation row. \\
    \bottomrule
  \end{tabularx}
\end{center}

Module-level helpers and constants:

\begin{pycode}
make_topology(topology, n_agents, *, edge_probability=0.3, rng=None,
              max_attempts=1000) -> np.ndarray       # adjacency of a static graph
proximity_adjacency(positions, comm_radius) -> np.ndarray   # disk graph
make_formation(shape, n_agents, radius=1.0) -> np.ndarray    # offsets d_i
algebraic_connectivity(adjacency) -> float   # second smallest eigenvalue of L
is_connected(adjacency) -> bool

TOPOLOGIES = ('ring', 'line', 'star', 'complete', 'erdos_renyi', 'proximity')
TASKS = ('consensus', 'formation')
DYNAMICS = ('single', 'double')
FORMATION_SHAPES = ('circle', 'line', 'grid', 'wedge')
\end{pycode}

\subsection{\texorpdfstring{\code{AJLATTEnv} and \code{AJLATTConfig}}{AJLATTEnv and AJLATTConfig}}

Module \pkg{env\_lib.ajlatt\_env}; Chapter~\ref{chap:ajlatt}.

\begin{pycode}
AJLATTEnv(config=None, render_mode=None, **overrides)
AJLATTConfig(map_name='obstacles04', num_robots=4, num_targets=1,
             max_episode_steps=120, sampling_period=0.5, ..., ci_solver='newton',
             target_policy='auto', ..., seed=None)   # dataclass, about 50 fields
TeamRewardWrapper(env)                # gymnasium.Wrapper
make(figID=0, *, render_mode=None, **kwargs) -> AJLATTEnv   # legacy factory
available_maps() -> list[str]
load_grid_map(map_name, margin2wall=0.5) -> GridMap
load_map(map_name) -> tuple[np.ndarray, list[float], list[float]]
\end{pycode}

\code{AJLATTEnv} is configured by an \code{AJLATTConfig} and/or keyword
overrides of its fields. \code{TeamRewardWrapper} returns the summed team
reward and a scalar \code{terminated} flag and keeps the per-robot values in
\code{info["agent\_rewards"]} and \code{info["agent\_terminated"]}.
\code{make} (also callable as \code{env\_lib.ajlatt\_env(...)}) accepts the
legacy parameter names with a \code{DeprecationWarning} and ignores unknown
names with a warning.

\begin{center}
  \small
  \begin{tabularx}{\textwidth}{@{}>{\raggedright\arraybackslash}p{6.9cm}>{\raggedright\arraybackslash}X@{}}
    \toprule
    Member & Description \\
    \midrule
    \multicolumn{2}{@{}l}{\emph{AJLATTEnv}} \\
    \code{single\_action\_space}, \code{single\_observation\_space} & Per-robot spaces; the joint spaces have a leading \code{num\_robots} axis. \\
    \code{observation\_layout()} & Dictionary of slices of the per-robot observation vector. \\
    \code{episode\_statistics()} & Dictionary of per-step traces of the current episode (\code{robot\_cov\_trace}, \code{target\_cov\_trace}, \code{robot\_error}, \code{target\_error}, \code{reward}), each \code{(steps + 1, num\_robots)}. \\
    \code{config}, \code{MAP} & The configuration and the loaded \code{GridMap}. \\
    \midrule
    \multicolumn{2}{@{}l}{\emph{AJLATTConfig}} \\
    \code{from\_kwargs(strict=True, **kwargs)} & Build from keyword arguments, translating legacy names. \\
    \code{add\_arguments(parser, prefix="")} & Register every scalar field as \code{-{}-<prefix><name>} on an \pkg{argparse} parser. \\
    \code{from\_namespace(namespace, prefix="", strict=False)} & Build from parsed arguments. \\
    \code{to\_dict()}, \code{from\_dict(data, strict=True)} & Plain-dictionary round trip. \\
    \code{replace(**changes)} & Validated copy with some fields replaced. \\
    \code{validate()}, \code{observation\_dim}, \code{field\_names()} & Range checks, per-robot observation size $7 + 4(n_R - 1) + 5$, field list. \\
    \bottomrule
  \end{tabularx}
\end{center}

\subsection{AJLATT maps}

Module \pkg{env\_lib.ajlatt\_env.maps}; Chapter~\ref{chap:ajlatt}.

\begin{pycode}
GridMap(map_path, margin2wall=0.5)
GridMap.from_array(occupancy, *, mapmin=(0.0, 0.0), mapres=0.4, margin2wall=0.5,
                   name='custom') -> GridMap
DynamicMap(map_dir_path, map_name, map_path=None, margin2wall=0.5)
DynamicMap.generate_map(chosen_idx=None, rot_angs=None, rng=None) -> np.ndarray
resolve_map_path(map_name) -> Path
MAP_DIR                                   # directory of the bundled maps
\end{pycode}

\begin{center}
  \small
  \begin{tabularx}{\textwidth}{@{}>{\raggedright\arraybackslash}p{6.9cm}>{\raggedright\arraybackslash}X@{}}
    \toprule
    \code{GridMap} method & Description \\
    \midrule
    \code{is\_blocked(start\_pos, end\_pos)} & Whether the straight line between two positions crosses an obstacle. \\
    \code{is\_blocked\_batch(start\_pos, end\_pos)} & Vectorised \code{is\_blocked} for arrays of shape \code{(n, 2)} or \code{(n, 3)}. \\
    \code{get\_closest\_obstacle(odom, ang\_res=0.05, fov=pi, r\_max=3.0, return\_pt=False)} & Range and bearing of the closest obstacle cell in a field of view. \\
    \code{get\_front\_obstacle(odom, r\_max=3.0)} & Range to the first obstacle straight ahead, or \code{None}. \\
    \code{is\_collision(pos, margin=None)} & Out of bounds or within \code{margin} of an occupied cell. \\
    \code{in\_bound(pos)} & At least \code{margin2wall} inside the map boundary. \\
    \code{se2\_to\_cell(pos)}, \code{cell\_to\_se2(cell)} & Conversions between positions and cell indices. \\
    \bottomrule
  \end{tabularx}
\end{center}

The submodules \pkg{maps.builder} (\code{build\_map}, \code{rasterize},
\code{load\_spec}, \code{save\_map}) and \pkg{maps.plotting}
(\code{plot\_map}) implement the console scripts \code{ajlatt-build-map} and
\code{ajlatt-plot-map}. The geometry helpers \code{bresenham2D},
\code{bresenham\_batch}, \code{se2\_to\_cell}, \code{cell\_to\_se2} and
\code{coord\_change2g} are exported as well.

\section{\texorpdfstring{\pkg{env\_lib.utils}}{env\_lib.utils}}\label{sec:api-utils}

\begin{pycode}
record_episode(env, policy=None, path=None, *, max_steps=None, seed=None,
               fps=None, render_every=1) -> list[np.ndarray]
save_animation(frames, path, fps=30.0, *, loop=0) -> Path
set_theme(theme)                          # "dark" (default) or "light"
get_theme(theme=None) -> Theme
register_theme(theme)
Theme(name, background, panel, grid, text, muted, accent, positive, negative,
      warning, categorical, sequential_cmap='viridis', diverging_cmap='RdBu_r')
DARK, LIGHT                               # the built-in themes
FigureWindow(render_mode, *, figsize=(8.0, 6.0), dpi=100, fps=30.0, title=None,
             theme=None)
MatplotlibRenderer(render_mode, *, figsize=(8.0, 6.0), dpi=100, fps=30.0,
                   title=None, theme=None, blit=True)
figure_to_rgb(fig) -> np.ndarray
style_axes(ax, theme, *, title=None, xlabel=None, ylabel=None, grid=True,
           show_spines=False, fontsize=9.0) -> Axes
with_alpha(color, alpha) -> tuple[float, float, float, float]
\end{pycode}

\code{record\_episode} and \code{save\_animation} are described in
Section~\ref{sec:examples-recording}. \code{Theme} is a dataclass of colours
with \code{color(index)} for the cycled categorical colours;
\code{set\_theme} affects renderers created afterwards. \code{FigureWindow}
owns one matplotlib figure (off-screen Agg canvas in \code{"rgb\_array"} mode,
a window in \code{"human"} mode) with \code{draw()}, \code{close()} and
\code{is\_open}. \code{MatplotlibRenderer} is the base class of the
environment renderers: subclasses implement \code{\_build(fig)} (create the
artists once) and \code{\_update(...)} (update their data), and users call
\code{render(...)}, \code{reset()} and \code{close()}.
\code{figure\_to\_rgb}, \code{style\_axes} and \code{with\_alpha} are helpers
for custom renderers.

\section{\texorpdfstring{\pkg{toolkit.plotkit}}{toolkit.plotkit}}\label{sec:api-plotkit}

Chapter~\ref{chap:plotkit}. All plot functions return the
\code{matplotlib.axes.Axes} they drew on.

\begin{pycode}
plot_shadow_curve(y, x=None, y_std=None, labels=None, colors=None, alpha=0.2,
                  ax=None, axis=0, figsize=(10, 6), title=None, xlabel=None,
                  ylabel=None, legend=True, grid=True, style='research',
                  legend_labels=None, x_tick_labels=None, y_tick_labels=None,
                  *, band='std', smoothing=None, **kwargs)
plot_learning_curves(runs, x=None, *, band='ci95', smoothing=None, colors=None,
                     alpha=0.2, ax=None, figsize=(10, 6), title=None,
                     xlabel='Step', ylabel='Return', legend=True, grid=True,
                     style='research', **kwargs)
plot_line(x, y, labels=None, colors=None, ax=None, figsize=(10, 6), title=None,
          xlabel=None, ylabel=None, legend=True, grid=True, style='research',
          **kwargs)
plot_scatter(x, y, labels=None, ...)       # same arguments as plot_line
plot_bar(x, height, labels=None, colors=None, ax=None, figsize=(10, 6),
         title=None, xlabel=None, ylabel=None, legend=True, grid=True,
         style='research', *, width=0.8, yerr=None, value_labels=False,
         horizontal=False, capsize=3.0, **kwargs)
plot_histogram(data, bins='auto', labels=None, colors=None, ax=None,
               figsize=(10, 6), title=None, xlabel=None, ylabel=None,
               legend=True, grid=True, style='research', *, density=False,
               alpha=0.45, histtype='stepfilled', **kwargs)
plot_heatmap(data, xlabels=None, ylabels=None, cmap='viridis', annot=False,
             ax=None, figsize=(8, 6), title=None, xlabel=None, ylabel=None,
             cbar=True, cbar_label=None, style='research', *, fmt=None,
             vmin=None, vmax=None, center=None, robust=False, mask=None,
             nan_color=None, aspect='auto', square=False, linewidths=None,
             linecolor='white', annot_kws=None, cbar_kws=None, **kwargs)
plot_gray_scale(data, xlabels=None, ylabels=None, ax=None, figsize=(8, 6),
                title=None, xlabel=None, ylabel=None, style='research',
                **kwargs)

save_figure(fig_or_ax, path, formats=None, dpi=300, transparent=False,
            **savefig_kwargs) -> list[Path]
style_context(style='research')           # context manager
set_research_style()
reset_style()
get_palette(name='research', n=None) -> list[str]

STYLE_PRESETS          # read-only; 'research', 'presentation', 'minimal', 'default'
RESEARCH_COLORS, RESEARCH_COLOR_LIST          # tab10, dict and list
OKABE_ITO_COLORS, OKABE_ITO_COLOR_LIST        # colour-blind-safe, dict and list
\end{pycode}

\code{python -m toolkit.plotkit} and the console script
\code{plotkit-gallery} run the gallery (Section~\ref{sec:plotkit-cli}). The
module \pkg{toolkit.plotkit.core} re-exports the same names for code written
against earlier versions.

\section{\texorpdfstring{\pkg{toolkit.neural\_toolkit}}{toolkit.neural\_toolkit}}\label{sec:api-neural}

Chapter~\ref{chap:neural}. Requires PyTorch. All networks are
\code{torch.nn.Module} subclasses; the shapes are given in
Table~\ref{tab:neural-shapes}.

\subsection{Layers}

\begin{pycode}
get_activation(name) -> nn.Module
build_mlp(input_dim, hidden_dims, output_dim=None, activation='relu',
          dropout=0.0, layer_norm=False, output_activation=None,
          bias=True) -> nn.Sequential
PositionalEncoding(d_model, dropout=0.1, max_len=5000)
ACTIVATION_NAMES   # ('elu', 'gelu', 'identity', 'leaky_relu', 'mish', 'relu',
                   #  'sigmoid', 'silu', 'softplus', 'swish', 'tanh')
\end{pycode}

\subsection{Policy and value networks}

\begin{pycode}
MLPPolicyNetwork(input_dim, output_dim, hidden_dims=(256, 256),
                 activation='relu', dropout=0.0, layer_norm=False, device='cpu')
CNPolicyNetwork(input_channels, output_dim, conv_dims=(32, 64, 128),
                fc_dims=(256, 256), kernel_sizes=None, strides=None,
                activation='relu', dropout=0.0, layer_norm=False, device='cpu')
RNPolicyNetwork(input_dim, output_dim, hidden_dim=256, num_layers=2,
                rnn_type='lstm', bidirectional=False, fc_dims=(256,),
                activation='relu', dropout=0.0, layer_norm=False, device='cpu')
TransformerPolicyNetwork(input_dim, output_dim, d_model=256, nhead=8,
                         num_layers=6, dim_feedforward=1024, dropout=0.1,
                         fc_dims=(256,), activation='relu', layer_norm=True,
                         device='cpu', max_len=5000)
CNNPolicyNetwork = CNPolicyNetwork; RNNPolicyNetwork = RNPolicyNetwork

MLPValueNetwork, CNNValueNetwork, RNNValueNetwork, TransformerValueNetwork
    # same arguments as the policy networks, with output_dim=1 by default
\end{pycode}

\subsection{Q-networks}

\begin{pycode}
MLPQNetwork(state_dim, action_dim, hidden_dims=(256, 256), activation='relu',
            dropout=0.0, layer_norm=False, device='cpu', continuous_action=False)
DuelingQNetwork(state_dim, action_dim, hidden_dims=(256, 256),
                activation='relu', dropout=0.0, layer_norm=False, device='cpu')
CNQNetwork(input_channels, action_dim, conv_dims=(32, 64, 128),
           fc_dims=(256, 256), kernel_sizes=None, strides=None,
           activation='relu', dropout=0.0, layer_norm=False, device='cpu')
RNQNetwork(state_dim, action_dim, hidden_dim=256, num_layers=2, rnn_type='lstm',
           bidirectional=False, fc_dims=(256,), activation='relu', dropout=0.0,
           layer_norm=False, device='cpu')
TransformerQNetwork(state_dim, action_dim, d_model=256, nhead=8, num_layers=6,
                    dim_feedforward=1024, dropout=0.1, fc_dims=(256,),
                    activation='relu', layer_norm=True, device='cpu',
                    max_len=5000)
CNNQNetwork = CNQNetwork; RNNQNetwork = RNQNetwork
# forward(state, action=None[, hidden or mask]): (B, A), or Q(s, a) of shape (B,)
\end{pycode}

\subsection{Encoders and decoders}

\begin{pycode}
MLPEncoder(input_dim, latent_dim, hidden_dims=(256, 256), activation='relu',
           dropout=0.0, layer_norm=False, device='cpu')
CNNEncoder(input_channels, latent_dim, conv_dims=(32, 64, 128, 256),
           fc_dims=(512, 256), kernel_sizes=None, strides=None,
           activation='relu', dropout=0.0, layer_norm=False, device='cpu',
           pooled_size=4)
RNNEncoder(input_dim, latent_dim, hidden_dim=256, num_layers=2, rnn_type='lstm',
           bidirectional=True, fc_dims=(256,), activation='relu', dropout=0.0,
           layer_norm=False, device='cpu')
TransformerEncoder(input_dim, latent_dim, d_model=256, nhead=8, num_layers=6,
                   dim_feedforward=1024, dropout=0.1, fc_dims=(256,),
                   activation='relu', layer_norm=True, device='cpu',
                   max_len=5000)
VariationalEncoder(input_dim, latent_dim, hidden_dims=(256, 256),
                   activation='relu', dropout=0.0, layer_norm=False,
                   device='cpu')          # forward -> (z, mu, logvar)
VariationalEncoder.kl_divergence(mu, logvar) -> Tensor   # static, shape (B,)

MLPDecoder(latent_dim, output_dim, hidden_dims=(256, 256), activation='relu',
           dropout=0.0, layer_norm=False, device='cpu')
CNNDecoder(latent_dim, output_channels, fc_dims=(512, 256),
           conv_dims=(256, 128, 64, 32), kernel_sizes=None, strides=None,
           initial_size=4, activation='relu', dropout=0.0, layer_norm=False,
           device='cpu', output_activation='sigmoid')   # property output_size
RNNDecoder(latent_dim, output_dim, hidden_dim=256, num_layers=2,
           rnn_type='lstm', max_seq_len=100, fc_dims=(256,), activation='relu',
           dropout=0.0, layer_norm=False, device='cpu')
TransformerDecoder(latent_dim, output_dim, d_model=256, nhead=8, num_layers=6,
                   dim_feedforward=1024, max_seq_len=100, dropout=0.1,
                   fc_dims=(256,), activation='relu', layer_norm=True,
                   device='cpu', max_len=None)
VariationalDecoder(latent_dim, output_dim, hidden_dims=(256, 256),
                   activation='relu', dropout=0.0, layer_norm=False,
                   device='cpu')
# RNNDecoder and TransformerDecoder: forward(z, seq_len=None)
\end{pycode}

\subsection{Factories}

\begin{pycode}
PolicyFactory.create_policy(policy_type, **kwargs)         # mlp cnn rnn transformer
ValueFactory.create_value_network(network_type, **kwargs)  # mlp cnn rnn transformer
QNetworkFactory.create_q_network(network_type, **kwargs)   # + dueling
EncoderFactory.create_encoder(encoder_type, **kwargs)      # + vae
DecoderFactory.create_decoder(decoder_type, **kwargs)      # + vae
\end{pycode}

The abstract base classes \code{BasePolicyNetwork}, \code{BaseValueNetwork},
\code{BaseQNetwork}, \code{BaseEncoder} and \code{BaseDecoder} are exported
for subclassing and type checks.

\subsection{\texorpdfstring{\code{NetworkUtils}}{NetworkUtils}}

All methods are static; Table~\ref{tab:neural-utils} describes them.

\begin{pycode}
initialize_weights(module, method='xavier_uniform', gain=1.0)
count_parameters(model) -> int
count_parameters_by_layer(model) -> dict[str, int]
freeze_layers(model, layer_names)
unfreeze_layers(model, layer_names)
get_grad_norm(model, norm_type=2.0) -> float
clip_grad_norm(model, max_norm, norm_type=2.0) -> float
create_mlp_layers(input_dim, hidden_dims, output_dim, activation='relu',
                  dropout=0.0, layer_norm=False, bias=True) -> nn.Sequential
create_conv_layers(input_channels, conv_dims, kernel_sizes=None, strides=None,
                   padding=None, activation='relu', dropout=0.0,
                   layer_norm=False) -> nn.Sequential
create_transposed_conv_layers(input_channels, conv_dims, kernel_sizes=None,
                              strides=None, padding=None, output_padding=None,
                              activation='relu', dropout=0.0,
                              layer_norm=False) -> nn.Sequential
create_attention_layer(input_dim, num_heads=8, dropout=0.1) -> nn.Module
create_positional_encoding(d_model, max_len=5000) -> nn.Parameter
get_activation(activation) -> nn.Module
apply_weight_decay(model, weight_decay, skip_list=None) -> list[dict]
create_optimizer(model, optimizer_type='adam', lr=0.001, weight_decay=0.0,
                 skip_list=None, **kwargs) -> torch.optim.Optimizer
create_scheduler(optimizer, scheduler_type='step', **kwargs)
save_checkpoint(model, optimizer, epoch, loss, filepath, **extra) -> Path
load_checkpoint(model, optimizer, filepath, map_location=None,
                weights_only=True, strict=True) -> tuple[int, float]
compute_output_size(input_size, kernel_size, stride=1, padding=0,
                    dilation=1) -> int
compute_transposed_output_size(input_size, kernel_size, stride=1, padding=0,
                               output_padding=0, dilation=1) -> int
compute_conv_output_size(input_size, conv_layers) -> tuple[int, int]
\end{pycode}

\subsection{Tabular tools}

\begin{pycode}
QTable(state_space_size, action_space_size=None, initial_value=0.0,
       dtype=np.float32, rng=None)
ValueTable(...same arguments...)
PolicyTable(...same arguments...)          # rows start uniform
BaseDiscreteTable                          # abstract base of the three tables
DiscreteEnvironment(state_space_size, action_space_size, rng=None)   # abstract

DiscreteTools.create_q_table(state_space_size, action_space_size,
                             initial_value=0.0, rng=None) -> QTable
DiscreteTools.create_value_table(state_space_size, initial_value=0.0,
                                 rng=None) -> ValueTable
DiscreteTools.create_policy_table(state_space_size, action_space_size,
                                  rng=None) -> PolicyTable
DiscreteTools.q_learning_update(q_table, state, action, reward, next_state,
                                gamma=0.99, alpha=0.1, done=False)
DiscreteTools.sarsa_update(q_table, state, action, reward, next_state,
                           next_action, gamma=0.99, alpha=0.1, done=False)
DiscreteTools.expected_sarsa_update(q_table, state, action, reward, next_state,
                                    policy_table=None, gamma=0.99, alpha=0.1,
                                    epsilon=0.0, done=False)
DiscreteTools.value_iteration_update(value_table, state, q_table, gamma=0.99)
DiscreteTools.policy_iteration_update(policy_table, state, q_table)
DiscreteTools.epsilon_greedy_policy(q_table, state, epsilon) -> int
DiscreteTools.softmax_policy(q_table, state, temperature) -> int
DiscreteTools.boltzmann_policy(q_table, state, temperature) -> int
DiscreteTools.ucb_policy(q_table, state, visit_counts,
                         exploration_constant=1.0) -> int
DiscreteTools.thompson_sampling_policy(q_table, state, visit_counts,
                                       prior_alpha=1.0, prior_beta=1.0) -> int
\end{pycode}

\section{\texorpdfstring{\pkg{toolkit.parakit}}{toolkit.parakit}}\label{sec:api-parakit}

Chapter~\ref{chap:toolkit}, Section~\ref{sec:parakit}.

\begin{pycode}
get_parameters(parser) -> dict[str, Any]
save_parameters(values, path_or_dir) -> Path
load_parameters(path) -> dict[str, Any]
apply_parameters(parser, values, strict=True,
                 validation_callbacks=None) -> argparse.ArgumentParser
parse_value(action, text, validator=None) -> Any
parse_bool(value) -> bool
format_value(action, value) -> str
describe_type(action) -> str
tunable_actions(parser) -> dict[str, argparse.Action]

ParameterTuner(parser, save_delay=30000, save_path=None, auto_save_enabled=True,
               validation_callbacks=None, inactivity_timeout=10)
ParameterTuner.tune() -> argparse.ArgumentParser
ParameterTuner.tune_parameters(parser, ...same arguments...)    # class method
ParameterAdjuster = ParameterTuner
TK_MISSING_MESSAGE                        # str: how to install Tk
\end{pycode}

\par\endgroup
